\section{Experiments}
\label{sec:experiments}

We evaluate \textsc{BHPOP} on two complementary benchmark datasets: proprietary cloud deployment workflows from Alibaba Cloud and open scientific workflow benchmarks from WFCommons. Our experiments investigate: (1) How does \textsc{BHPOP} perform relative to state-of-the-art process mining methods? (2) How does inference quality depend on the critical pair coverage of observed traces? (3) What is the effect of the likelihood smoothing parameter $\varepsilon$?

\subsection{Experimental Setup}

\paragraph{Dataset 1: Alibaba Cloud Deployment Workflows.}
We use six cloud infrastructure deployment scenarios from Alibaba Cloud, representing real-world DevOps workflows for provisioning compute, networking, and database resources. Table~\ref{tab:scenarios} summarizes the scenarios. Each scenario has a ground-truth partial order (DAG) specifying valid task dependencies, with 5--12 tasks and 5--14 edges. We collected execution traces from both expert demonstrations and automated agents, filtering for successful deployments and removing unknown actions, yielding a combined dataset of human and machine-generated traces.

\paragraph{Dataset 2: WFCommons Scientific Workflows.}
To validate our method on \emph{open, reproducible benchmarks} with larger-scale DAGs and richer structural complexity, we use the WFCommons scientific workflow repository~\citep{coleman2022wfcommons}---a community standard for workflow research. We focus on the \textbf{Montage} astronomical image mosaicking workflow~\citep{rynge2013producing}, which creates composite images from multiple telescope observations. Montage exhibits characteristic \emph{fan-out} (parallel projection) and \emph{fan-in} (aggregation) patterns that stress-test structural recovery algorithms.

We use two Montage instances as \emph{hierarchical assessors}:
\begin{itemize}[leftmargin=*,nosep]
    \item \texttt{montage-2mass-005d}: 58 tasks, 114 edges (0.05° sky region)
    \item \texttt{montage-2mass-01d}: 103 tasks, 231 edges (0.1° sky region)
\end{itemize}
These instances share 16 common tasks while differing in scale, enabling hierarchical inference where shared task dependencies are pooled across instances. Unlike the Alibaba Cloud scenarios where traces come from sequential agent interactions, WFCommons provides execution logs from parallel schedulers, testing \textsc{BHPOP}'s ability to recover partial orders from highly parallel traces. See Appendix~\ref{app:wfcommons} for full details.

\begin{table}[t]
\caption{Cloud deployment scenario characteristics.}
\label{tab:scenarios}
\centering
\small
\begin{tabular}{lcc}
\toprule
\textbf{Scenario} & \textbf{Tasks} & \textbf{Edges} \\
\midrule
\texttt{simple\_ecs} & 5 & 5 \\
\texttt{dual\_zone\_ecs\_slb} & 7 & 8 \\
\texttt{eip\_slb\_ecs} & 9 & 10 \\
\texttt{slb\_ecs\_redis} & 9 & 10 \\
\texttt{dual\_zone\_ecs\_slb\_rds} & 10 & 12 \\
\texttt{slb\_ecs\_rds} & 12 & 14 \\
\bottomrule
\end{tabular}
\end{table}

\paragraph{Trace Augmentation.}
To systematically study the effect of trace coverage, we augment real traces with synthetic linear extensions sampled uniformly from the ground-truth DAG using Kahn's algorithm with random frontier selection. We vary the \emph{critical pair coverage} (CP-Cov)---the fraction of incomparable pairs observed in both orderings---from 0.5 to 1.0.

\paragraph{Baselines.}
We compare against three baselines:
\begin{itemize}[leftmargin=*,nosep]
    \item \textbf{Majority}: For each pair $(i,j)$, add edge $i \to j$ if $i$ precedes $j$ in $>50\%$ of co-occurrences. Cycles are broken greedily by removing the weakest edge.
    \item \textbf{Inductive Miner} \citep{leemans2013discovering}: A state-of-the-art process mining algorithm that discovers process trees via recursive log splitting. We use the IMf variant with noise threshold 0.
    \item \textbf{Heuristics Miner} \citep{weijters2006process}: A frequency-based process mining algorithm that uses dependency measures to discover precedence relations, with dependency threshold 0.5.
\end{itemize}

\paragraph{Evaluation Metrics.}
We evaluate recovered partial orders using:
\begin{itemize}[leftmargin=*,nosep]
    \item \textbf{F1 Score}: Harmonic mean of precision and recall on edge recovery, comparing inferred cover relations to ground truth.
    \item \textbf{Structural Hamming Distance (SHD)}: Number of edge additions, deletions, and reversals needed to transform the inferred graph to ground truth.
    \item \textbf{Feasibility}: Fraction of observed traces that are valid linear extensions of the inferred partial order.
\end{itemize}

\paragraph{Implementation Details.}
We run MCMC inference for $10^6$ iterations with 50\% burn-in. The posterior is aggregated by thresholding the mean at 0.5, followed by transitive reduction. We sweep the likelihood smoothing parameter $\varepsilon \in \{0.001, 0.005, 0.01, 0.05, 0.1\}$ and report results across all 35 experimental conditions (7 CP-Cov levels $\times$ 5 $\varepsilon$ values) per scenario.

\subsection{Main Results}

Table~\ref{tab:main_results} presents the main comparison. \textsc{BHPOP} achieves the highest mean F1 score (0.652) and is the only method capable of achieving perfect reconstruction (F1 = 1.0) on multiple scenarios. While the Majority baseline achieves comparable mean performance, it exhibits much higher variance and lower feasibility.

\begin{table}[t]
\caption{Performance comparison across all scenarios and experimental conditions. Best results in \textbf{bold}.}
\label{tab:main_results}
\centering
\small
\begin{tabular}{lccccc}
\toprule
\textbf{Method} & \textbf{F1}$\uparrow$ & \textbf{F1 (max)}$\uparrow$ & \textbf{SHD}$\downarrow$ & \textbf{Feas.}$\uparrow$ \\
\midrule
\textsc{BHPOP} (Ours) & \textbf{0.652} $\pm$ 0.21 & \textbf{1.000} & \textbf{8.9} $\pm$ 6.7 & 0.78 $\pm$ 0.26 \\
Majority & 0.577 $\pm$ 0.31 & 1.000 & 8.9 $\pm$ 6.9 & 0.48 $\pm$ 0.48 \\
Heuristics Miner & 0.574 $\pm$ 0.10 & 0.737 & 8.0 $\pm$ 3.2 & 0.18 $\pm$ 0.25 \\
Inductive Miner & 0.463 $\pm$ 0.18 & 0.727 & 10.6 $\pm$ 5.5 & \textbf{1.00} $\pm$ 0.00 \\
\bottomrule
\end{tabular}
\end{table}

\paragraph{Per-Scenario Analysis.}
Figure~\ref{fig:per_scenario} shows F1 scores by scenario. \textsc{BHPOP} achieves perfect reconstruction (F1 = 1.0) on three scenarios: \texttt{simple\_ecs}, \texttt{dual\_zone\_ecs\_slb}, and \texttt{dual\_zone\_ecs\_slb\_rds}. The method performs best on scenarios with simpler dependency structures (5--10 tasks), while more complex scenarios (\texttt{slb\_ecs\_rds} with 12 tasks) remain challenging for all methods.

\begin{table}[t]
\caption{\textsc{BHPOP} performance by scenario (best configuration).}
\label{tab:per_scenario}
\centering
\small
\begin{tabular}{lcccc}
\toprule
\textbf{Scenario} & \textbf{Avg F1} & \textbf{Best F1} & \textbf{Best $\varepsilon$} \\
\midrule
\texttt{simple\_ecs} & 0.852 & 1.000 & 0.05--0.1 \\
\texttt{dual\_zone\_ecs\_slb} & 0.830 & 1.000 & 0.01--0.1 \\
\texttt{dual\_zone\_ecs\_slb\_rds} & 0.656 & 1.000 & 0.1 \\
\texttt{slb\_ecs\_redis} & 0.610 & 0.900 & 0.05--0.1 \\
\texttt{eip\_slb\_ecs} & 0.489 & 0.842 & 0.1 \\
\texttt{slb\_ecs\_rds} & 0.478 & 0.875 & 0.1 \\
\bottomrule
\end{tabular}
\end{table}

\subsection{Effect of Critical Pair Coverage}

Figure~\ref{fig:cpcov_effect} shows the relationship between CP-Cov and F1 score. As expected, higher coverage leads to better reconstruction quality. At CP-Cov = 1.0 (full coverage of incomparable pairs), \textsc{BHPOP} achieves mean F1 of 0.83 with peak performance of 0.92. Notably, even at CP-Cov = 0.5, \textsc{BHPOP} maintains reasonable performance (mean F1 $\approx$ 0.55), demonstrating robustness to limited trace diversity.

\subsection{Effect of Smoothing Parameter $\varepsilon$}

The likelihood smoothing parameter $\varepsilon$ has a pronounced effect on inference quality (Figure~\ref{fig:epsilon_effect}). Counterintuitively, larger values of $\varepsilon$ (0.05--0.1) consistently outperform smaller values (0.001--0.01). We attribute this to two factors:

\begin{enumerate}[leftmargin=*,nosep]
    \item \textbf{MCMC mixing}: Small $\varepsilon$ creates a highly peaked likelihood, causing the Markov chain to get trapped in local modes. Larger $\varepsilon$ smooths the posterior landscape, enabling better exploration.
    \item \textbf{Regularization}: The smoothing acts as implicit regularization against overfitting to trace-specific orderings, improving generalization to the true partial order.
\end{enumerate}

This finding suggests practitioners should use $\varepsilon \in [0.05, 0.1]$ for robust inference. The optimal trade-off occurs around $\varepsilon = 0.05$--$0.1$, where the likelihood remains informative while permitting adequate posterior exploration.

\subsection{Computational Efficiency and Scalability}
\label{sec:scalability}

We report detailed runtime and memory measurements to address practical deployment concerns.

\paragraph{Runtime Analysis.}
Table~\ref{tab:runtime} summarizes the computational cost. Each MCMC run ($10^6$ iterations) completes in approximately \textbf{9 minutes} on a single core of an Apple M1 CPU. With 8 parallel workers, the complete experiment suite (35 configurations $\times$ 6 scenarios = 210 runs) finishes in \textbf{5.3 hours wall-clock time}, corresponding to an effective per-run time of \textbf{1.1 minutes}.

\begin{table}[htbp]
\caption{Computational cost of BHPOP inference.}
\label{tab:runtime}
\centering
\small
\begin{tabular}{lc}
\toprule
\textbf{Metric} & \textbf{Value} \\
\midrule
MCMC iterations & $10^6$ \\
Single-core time per run & 9 min \\
Parallel workers & 8 \\
Wall-clock time (35 runs) & 5.3 hours \\
Effective time per run & 1.1 min \\
Peak memory per run & $<$500 MB \\
Posterior storage (H\_trace) & 9.6 MB \\
\bottomrule
\end{tabular}
\end{table}

\paragraph{Scaling Behavior.}
The per-iteration complexity is $O(n^2 \cdot m)$ where $n$ is the number of tasks and $m$ is the number of traces, dominated by likelihood evaluation over the $n \times n$ adjacency matrix. For our scenarios (5--12 tasks, 10--100 traces), this remains tractable. Scaling to larger problems can be addressed via:
\begin{itemize}[leftmargin=*,nosep]
    \item \textbf{Subsampling}: Using CP-Cov-guided trace selection (Section~5.3) to reduce $m$ while preserving information content.
    \item \textbf{Parallelization}: Independent MCMC chains across assessors achieve near-linear speedup.
    \item \textbf{Early stopping}: Convergence diagnostics (Gelman-Rubin $\hat{R}$) can reduce iterations by 50--80\% for well-specified models.
\end{itemize}

\paragraph{Practical Deployment.}
For typical deployment scenarios:
\begin{itemize}[leftmargin=*,nosep]
    \item \textbf{Online inference} (single workflow, 10 tasks, 50 traces): $<$2 minutes with $10^5$ iterations.
    \item \textbf{Batch analysis} (6 workflows, 100 traces each): $<$1 hour on 8 cores.
    \item \textbf{Hyperparameter tuning}: Grid search over $\varepsilon$ adds linear overhead; we recommend $\varepsilon \in [0.05, 0.1]$ as a robust default.
\end{itemize}

Compared to process mining baselines (Inductive/Heuristics Miner), which run in seconds, BHPOP trades speed for principled uncertainty quantification and superior peak accuracy. This trade-off is justified for applications requiring high-fidelity structural recovery, such as compliance verification and workflow optimization.

\subsection{Unique Benefits of Bayesian Posterior Inference}
\label{sec:bayesian_benefits}

We highlight three key advantages of our likelihood-based Bayesian approach over point-estimate baselines.

\paragraph{(1) Full Posterior Uncertainty Quantification.}
Unlike process mining methods that output a single DAG, \textsc{BHPOP} provides a \emph{complete posterior distribution} over partial orders. Figure~\ref{fig:exp35_posterior} shows posterior parameter distributions from a representative experiment (Exp~35: CP-Cov=1.0, $\varepsilon$=0.1). The inferred parameters exhibit well-behaved posteriors:

\begin{table}[htbp]
\caption{Posterior parameter estimates (Exp~35, 500K post-burn-in samples).}
\label{tab:posterior_params}
\centering
\small
\begin{tabular}{lccc}
\toprule
\textbf{Parameter} & \textbf{Mean} & \textbf{95\% CI} & \textbf{Interpretation} \\
\midrule
$\rho$ (edge prior) & 0.957 & [0.924, 0.977] & Strong edge preference \\
$\tau$ (assessor var.) & 0.428 & [0.036, 0.832] & Moderate heterogeneity \\
$K$ (latent dimension) & 6.93 & [3, 11] & ~7 latent factors \\
$\beta$ (softmax temp.) & 1.13 & [0.18, 2.72] & Moderate sharpness \\
\midrule
Log-likelihood & -239.2 & [$-$308, $-$194] & Well-converged \\
\bottomrule
\end{tabular}
\end{table}

This enables practitioners to quantify confidence in individual edges, identify uncertain dependencies, and propagate uncertainty to downstream decisions.

\paragraph{(2) Robustness via Likelihood Smoothing.}
The $\varepsilon$-smoothed likelihood (Eq.~X) provides principled robustness to trace noise:
\begin{itemize}[leftmargin=*,nosep]
    \item \textbf{Small $\varepsilon$}: Sharp likelihood, high sensitivity to violations, poor MCMC mixing.
    \item \textbf{Large $\varepsilon$}: Smooth likelihood, robust to noise, better exploration.
\end{itemize}
Our experiments (Section~5.4) show $\varepsilon \in [0.05, 0.1]$ achieves optimal balance. This \emph{learned} robustness differs fundamentally from ad-hoc filtering in process mining methods.

\paragraph{(3) Recoverability Guarantees via CP-Cov.}
The critical pair coverage metric provides a \emph{theoretical certificate} for recovery:
\begin{itemize}[leftmargin=*,nosep]
    \item At CP-Cov = 1.0: All incomparable pairs observed in both orderings $\Rightarrow$ full identifiability.
    \item At CP-Cov $< 1.0$: Some pairs unobserved $\Rightarrow$ partial identifiability with quantified uncertainty.
\end{itemize}
Figure~\ref{fig:cpcov_effect} confirms this theoretical prediction empirically: F1 increases monotonically with CP-Cov, providing actionable guidance for trace collection.

\subsection{Discussion}

Our experiments reveal several insights:

\textbf{(1) Probabilistic modeling outperforms heuristics.} While process mining methods like Heuristics Miner achieve consistent but limited performance, \textsc{BHPOP}'s probabilistic framework enables principled uncertainty quantification and achieves superior peak performance.

\textbf{(2) Trace diversity matters more than quantity.} The critical pair coverage metric effectively predicts reconstruction quality. Practitioners should prioritize collecting traces that exercise different valid orderings rather than many similar traces.

\textbf{(3) Likelihood smoothing is crucial.} The $\varepsilon$ parameter should be tuned to balance likelihood informativeness against MCMC mixing. Our results suggest $\varepsilon \approx 0.05$--$0.1$ as a robust default.

\textbf{Limitations.} Our evaluation is limited to six cloud deployment scenarios with 5--12 tasks. Scalability to larger partial orders (50+ tasks) remains to be validated. Additionally, the synthetic trace augmentation may not fully capture the diversity of real-world execution traces.

%% ============================================
%% APPENDIX
%% ============================================

\appendix

\section{WFCommons Workflow Details}
\label{app:wfcommons}

\subsection{Dataset Availability and Reproducibility}

All WFCommons workflow instances are publicly available under open licenses:
\begin{itemize}[leftmargin=*,nosep]
    \item \textbf{Repository}: \url{https://github.com/wfcommons/WfInstances}
    \item \textbf{Montage documentation}: \url{https://github.com/wfcommons/WfInstances/blob/main/pegasus/montage/README.md}
\end{itemize}

\subsection{Montage Workflow Structure}

\textbf{Montage} is a portable toolkit for assembling astronomical images into composite mosaics~\citep{rynge2013producing}. The workflow processes multiple input images through several stages, as illustrated in Figure~\ref{fig:montage_dag}.

\begin{figure}[htbp]
    \centering
    \includegraphics[width=0.9\linewidth]{montage.png}
    \caption{\textbf{Ground-truth structure of the Montage scientific workflow.} This DAG represents a complex astronomical image mosaicking task from the WFCommons benchmark~\citep{rynge2013producing}. Unlike linear agent traces, Montage exhibits high structural parallelism with characteristic ``fan-out'' (parallel processing) and ``fan-in'' (aggregation) patterns, serving as a rigorous stress test for our structural recovery algorithm on data-intensive topologies.}
    \label{fig:montage_dag}
\end{figure}

Table~\ref{tab:montage_tasks} describes the task types in Montage workflows:

\begin{table}[htbp]
\caption{Montage workflow task types and their functions.}
\label{tab:montage_tasks}
\centering
\small
\begin{tabular}{llc}
\toprule
\textbf{Task Type} & \textbf{Function} & \textbf{Count} \\
\midrule
\texttt{mProject} & Reproject images to common coordinates & $n$ \\
\texttt{mDiffFit} & Compute pairwise image differences & $O(n^2)$ \\
\texttt{mConcatFit} & Concatenate difference fits & 1--3 \\
\texttt{mBgModel} & Compute background correction model & 1--3 \\
\texttt{mBackground} & Apply background corrections & $n$ \\
\texttt{mImgtbl} & Generate metadata table & 1--3 \\
\texttt{mAdd} & Co-add images into final mosaic & 1--3 \\
\texttt{mViewer} & Generate visualization outputs & 1--4 \\
\bottomrule
\end{tabular}
\end{table}

\subsection{Hierarchical Assessor Structure}

The two Montage instances form a hierarchical partial order where:

\begin{enumerate}[leftmargin=*,nosep]
    \item \textbf{Shared tasks (16):} Both instances contain tasks with identical IDs for common operations (e.g., \texttt{mProject\_ID0000001} through \texttt{mProject\_ID0000004}).
    \item \textbf{Instance-specific tasks:} \texttt{montage-005d} has 42 unique tasks; \texttt{montage-01d} has 87 unique tasks.
    \item \textbf{Global task space:} 145 unique tasks total ($58 + 103 - 16$ shared).
\end{enumerate}

This structure enables \emph{hierarchical inference} where information about shared tasks is pooled across instances, while instance-specific dependencies are learned separately.

\subsection{Trace Generation Protocol}

\textbf{Execution-derived linearizations} are generated by:
\begin{enumerate}[leftmargin=*,nosep]
    \item Running the workflow under the Pegasus workflow management system
    \item Recording task start/completion timestamps from execution logs
    \item Ordering tasks by start time to obtain a total order
    \item Augmenting with synthetic traces via Kahn's algorithm (uniform random topological sorts)
\end{enumerate}

\subsection{Dataset Comparison}

Table~\ref{tab:dataset_comparison} summarizes the key differences between our two evaluation datasets:

\begin{table}[htbp]
\caption{Comparison of Alibaba Cloud and WFCommons datasets.}
\label{tab:dataset_comparison}
\centering
\small
\begin{tabular}{lcc}
\toprule
\textbf{Aspect} & \textbf{Alibaba Cloud} & \textbf{WFCommons} \\
\midrule
Domain & Cloud deployment & Scientific computing \\
Trace source & Sequential agent & Parallel scheduler \\
Scale & 5--12 tasks & 58--619 tasks \\
Structure & Sequential chains & Fan-out/fan-in \\
Ground truth & Expert-defined & Workflow spec \\
Hierarchy & Independent & Shared tasks \\
Availability & Proprietary & Open benchmark \\
\bottomrule
\end{tabular}
\end{table}

The complementary nature of these datasets validates \textsc{BHPOP}'s generality across both real-world interactive traces with noise and standardized benchmarks with controlled parallelism.
